% oss-plantilla: oss-2026.11
% Preámbulo corporativo fijo. No editar en sd-NN.tex: exportar_latex.py verificar lo compara byte a byte.
% Para cambiar el formato, editar esta plantilla u oss.sty y subir la versión de la primera línea.
\PassOptionsToPackage{unicode}{hyperref}
\PassOptionsToPackage{hyphens}{url}
\documentclass[11pt]{article}
\usepackage{amsmath,amssymb}
\usepackage{iftex}
\ifPDFTeX
  % Prism y otros editores en la nube: pdfLaTeX con Helvetica como familia por defecto.
  \usepackage[T1]{fontenc}
  \usepackage[utf8]{inputenc}
  \usepackage{textcomp}
  \usepackage[scaled=0.92]{helvet}
  \renewcommand{\familydefault}{\sfdefault}
  \DeclareUnicodeCharacter{2264}{\ensuremath{\leq}}
  \DeclareUnicodeCharacter{2265}{\ensuremath{\geq}}
  \DeclareUnicodeCharacter{2260}{\ensuremath{\neq}}
  \DeclareUnicodeCharacter{2248}{\ensuremath{\approx}}
  \DeclareUnicodeCharacter{2190}{\ensuremath{\leftarrow}}
  \DeclareUnicodeCharacter{2192}{\ensuremath{\rightarrow}}
  \DeclareUnicodeCharacter{25CF}{\ensuremath{\bullet}}
  \DeclareUnicodeCharacter{2713}{\ensuremath{\checkmark}}
\else
  % XeLaTeX/LuaLaTeX (compilación oficial local): DejaVu Sans.
  \usepackage{fontspec}
  \defaultfontfeatures{Ligatures=TeX,Scale=MatchLowercase}
  \setmainfont{DejaVu Sans}
  \setsansfont{DejaVu Sans}
  \setmonofont{DejaVu Sans Mono}
\fi
\usepackage[letterpaper,margin=20mm]{geometry}
\usepackage{microtype}
\usepackage{parskip}
\usepackage{xcolor}
\usepackage{longtable,booktabs,array}
\usepackage{calc}
\usepackage{etoolbox}
\makeatletter
\patchcmd\longtable{\par}{\if@noskipsec\mbox{}\fi\par}{}{}
\makeatother
\usepackage{footnotehyper}
\makesavenoteenv{longtable}
\usepackage{graphicx}
\makeatletter
\newsavebox\pandoc@box
\newcommand*\pandocbounded[1]{%
  \sbox\pandoc@box{#1}%
  \Gscale@div\@tempa{\textheight}{\dimexpr\ht\pandoc@box+\dp\pandoc@box\relax}%
  \Gscale@div\@tempb{\linewidth}{\wd\pandoc@box}%
  \ifdim\@tempb\p@<\@tempa\p@\let\@tempa\@tempb\fi
  \ifdim\@tempa\p@<\p@\scalebox{\@tempa}{\usebox\pandoc@box}%
  \else\usebox{\pandoc@box}%
  \fi}
\def\maxwidth{\ifdim\Gin@nat@width>\linewidth\linewidth\else\Gin@nat@width\fi}
\def\maxheight{\ifdim\Gin@nat@height>\textheight\textheight\else\Gin@nat@height\fi}
\makeatother
\setkeys{Gin}{width=\maxwidth,height=\maxheight,keepaspectratio}
\makeatletter
\def\fps@figure{htbp}
\makeatother
\usepackage[normalem]{ulem}
\providecommand{\st}[1]{\sout{#1}}
\usepackage{fancyvrb}
\setlength{\emergencystretch}{3em}
\providecommand{\tightlist}{%
  \setlength{\itemsep}{0pt}\setlength{\parskip}{0pt}}
\providecommand{\passthrough}[1]{#1}
\setcounter{secnumdepth}{-\maxdimen}
\ifPDFTeX
  \usepackage[spanish]{babel}
\else
  \usepackage[bidi=default]{babel}
  \babelprovide[main,import]{spanish}
\fi
\let\LanguageShortHands\languageshorthands
\def\languageshorthands#1{}
\usepackage{bookmark}
\usepackage{xurl}
\urlstyle{same}
\hypersetup{pdflang={es-CL},pdfcreator={Only Simple Solutions - exportar_latex.py}}
\usepackage{oss}
\begin{document}
\renewcommand{\tablename}{Tabla}

\ossCover{Presentación de la empresa}{Subdocumento 1}
\tableofcontents
\clearpage
\renewcommand{\listfigurename}{Lista de figuras}
\listoffigures
\clearpage
\renewcommand{\listtablename}{Lista de tablas}
\listoftables
\clearpage

\section{Capítulo 1 · Introducción}\label{capuxedtulo-1-introducciuxf3n}

Este capítulo presenta los antecedentes institucionales de Only Simple Solutions y explica su relación con el contrato de Multitiendas Ancoa S.A. Describe la trayectoria de la empresa, sus capacidades instaladas y los productos y servicios que ofrece; su estructura organizacional y dotación y las políticas, instancias y responsables de su gobierno interno de calidad, seguridad de la información y gestión del conocimiento. También resume la experiencia pertinente y las certificaciones institucionales, explica la organización prevista para abordar el proyecto y presenta las alianzas tecnológicas de la empresa.

La pertinencia de estas capacidades se relaciona con el problema y el alcance desarrollados en los capítulos 2 y 3. En Ancoa, la integración de las plataformas existentes debe preservar la continuidad de la operación comercial y respetar la separación de los datos del negocio de retail y del emisor de crédito fiscalizado. Los capítulos 4 y 5 desarrollan la respuesta de arquitectura y gestión de datos a estas condiciones. A su vez, el gobierno interno presentado aquí proporciona el contexto organizacional para los planes de trabajo y calidad y los servicios de operación descritos en los capítulos 7, 9, 10 y 11.

El detalle de los proyectos de referencia se presenta en el Formulario T-6, entregado por separado, mientras que el anexo A.1 complementa la presentación institucional con la referencia al sitio web corporativo. Las certificaciones y alianzas declaradas requieren sus respectivos documentos de respaldo. La estructura general para ejecutar el contrato se resume en la sección 1.5, el equipo nominado, sus credenciales y dedicaciones se desarrollarán en el capítulo 12 y el Formulario T-8. Las alianzas específicas para este proyecto se detallan en la sección 12.3.

\section{1.1. Presentación de la empresa}\label{presentaciuxf3n-de-la-empresa}

Only Simple Solutions es una compañía integradora de soluciones tecnológicas y desarrollo de software a medida, fundada en Santiago el año 2016 para atender los desafíos más exigentes del mercado corporativo latinoamericano. Nuestra línea de negocio se centra exclusivamente en la ingeniería, orquestación y modernización de sistemas transaccionales de misión crítica. A lo largo de nuestra historia, nos hemos especializado en resolver la complejidad técnica de grandes corporaciones que operan en industrias de transaccionalidad masiva y fuertemente reguladas, como el sector bancario, la logística a gran escala, telecomunicaciones y retail.

Entendemos que en estos entornos corporativos, la tecnología no debe enfocarse exclusivamente en respaldar y agilizar la operación del negocio, sino que también garantizar, de forma trazable y auditable, el cumplimiento de las normas aplicables. Por ello, nuestro enfoque va más allá de soluciones genéricas, nos especializamos en el diseño y construcción de software a medida orientados a modernizar plataformas legadas y evolucionar arquitecturas monolíticas hacia ecosistemas distribuidos, escalables y seguros. Todo ello bajo un enfoque de transición controlada, que permite minimizar riesgos y preservar la continuidad de los servicios durante el proceso de transformación.

Contamos con 85 profesionales distribuidos en dos sedes, la Casa Matriz en Santiago y la Oficina Regional de Concepción, y con una trayectoria de 14 proyectos desde 2016, de los que 6 se mantienen vigentes con clientes de banca, retail, logística, telecomunicaciones y seguros.

Nuestra visión es ser, hacia 2032, la compañía de referencia de Latinoamérica en la modernización de sistemas transaccionales de misión crítica, y una organización en la que los profesionales de la región quieran crecer. Para alcanzarla nos fijamos tres metas de largo plazo:

\begin{itemize}
\tightlist
\item
  Expansión internacional: Llevar nuestra operación más allá de Chile hacia los principales mercados de la región como Perú, Colombia, Argentina y México, sostenida en sociedades locales, alianzas y equipos en terreno.
\item
  Activo propio y transferible: Convertir la experiencia acumulada en plataformas, componentes y metodologías reutilizables y certificables, en lugar de soluciones que se rehacen para cada cliente.
\item
  Cumplimiento como ventaja competitiva: Lograr que el buen desempeño comercial y el cumplimiento normativo dejen de verse como cosas en tensión, y sostenerlo en el tiempo como una ventaja verificable para nuestros clientes.
\end{itemize}

Nuestros valores se sustentan en la responsabilidad y el compromiso continuo, entendiendo que los sistemas críticos de nuestros clientes no pueden fallar en ningún caso. Trabajamos siempre guiados por la probidad corporativa, la protección integral de los derechos de los consumidores y el respeto absoluto por las restricciones legales de cada rubro. Por último, creemos en la transparencia y en el sentido práctico, creando soluciones tecnológicas que dejen un rastro claro y auditable, pero que al mismo tiempo faciliten el trabajo diario de las personas en lugar de complicarlo.

\subsection{1.1.1. Capacidades instaladas}\label{capacidades-instaladas}

Para garantizar que los sistemas críticos operen sin interrupciones y cumplan con todas las exigencias normativas, hemos estructurado nuestras capacidades operativas con un objetivo claro: neutralizar los riesgos técnicos y operativos antes de que impacten al negocio.

En lugar de disponer de áreas aisladas, articulamos nuestro talento en equipos multidisciplinarios que se complementan para asegurar la escalabilidad, la seguridad y la continuidad de cada solución:

\begin{itemize}
\tightlist
\item
  Fábrica de software (Construcción y Evolución): Más allá de escribir código, diseñamos e integramos soluciones informáticas pensadas para durar y adaptarse al cambio. Al hacernos cargo de la transición de sistemas legados y la modernización de arquitecturas complejas, evitamos los cuellos de botella y la obsolescencia tecnológica. Así, garantizamos que las plataformas evolucionen de manera controlada y sin poner en riesgo la operación diaria.
\item
  Célula de datos (Confiabilidad y Gobierno): En escenarios de alta transaccionalidad, la inconsistencia o pérdida de información no es una opción. Nuestro equipo especializado asume el control del gobierno, la integración y el procesamiento de grandes volúmenes de información. Con esto, eliminamos la incertidumbre en los procesos de migración y aseguramos que el cliente siempre cuente con datos íntegros y exactos para la toma de decisiones.
\item
  Célula de ciberseguridad (Protección Integral): Entendemos que la seguridad debe proteger el negocio, no entorpecerlo. Este equipo se encarga de blindar la infraestructura y las aplicaciones desde el diseño mismo de la solución, implementando controles de acceso, segmentación de redes y estrategias de mitigación de riesgos que se adaptan a la realidad de la organización, protegiendo sus activos críticos sin sacrificar la agilidad en el trabajo diario.
\item
  Mesa de servicio (Continuidad Operacional): Para que la tecnología cumpla su propósito de facilitar el día a día, acompañamos la operación con un soporte multinivel experto. Nos adelantamos a los problemas mediante el monitoreo constante y gestionamos cualquier incidente de forma ágil. Esta capacidad nos permite aislar las fallas técnicas, coordinar los escalamientos necesarios y asegurar que los usuarios finales siempre reciban un servicio estable y de alta calidad.
\end{itemize}

\subsection{1.1.2. Líneas de negocio, productos y servicios ofrecidos}\label{luxedneas-de-negocio-productos-y-servicios-ofrecidos}

Only Simple Solutions estructura su oferta en torno a la modernización, integración y operación de soluciones tecnológicas para organizaciones que requieren alta disponibilidad, procesamiento transaccional, seguridad y trazabilidad. La compañía se orienta principalmente al desarrollo de soluciones a medida, adaptadas a las características operativas, regulatorias y tecnológicas de cada cliente.

Nuestras principales líneas de negocio son las siguientes:

\begin{itemize}
\tightlist
\item
  Modernización de sistemas legados: Evolución de plataformas monolíticas y aplicaciones heredadas hacia arquitecturas distribuidas, modulares, escalables y mantenibles, procurando reducir riesgos de transición y mantener la continuidad operacional durante los procesos de cambio.
\item
  Desarrollo de software a medida: Diseño, construcción y evolución de aplicaciones transaccionales y plataformas digitales ajustadas a los procesos específicos de cada organización, incluyendo soluciones para retail, servicios financieros, logística y otros entornos de alta transaccionalidad.
\item
  Arquitectura e integración de sistemas: Diseño de arquitecturas on-premise, cloud e híbridas, junto con la integración de aplicaciones, plataformas y servicios mediante APIs, servicios de integración y mecanismos de interoperabilidad que permitan disminuir la dependencia de integraciones punto a punto.
\item
  Operación y soporte multinivel: Prestación de servicios de soporte y operación con acuerdos de nivel de servicio, monitoreo, gestión de incidentes y continuidad del negocio.
\item
  Aseguramiento de calidad y pruebas: Definición y ejecución de estrategias de pruebas funcionales, de integración, carga, resiliencia y aceptación de usuario, con el propósito de reducir riesgos antes de la puesta en producción de soluciones críticas.
\end{itemize}

No basamos su propuesta de valor en un catálogo cerrado de productos de software estándar, nuestro modelo se orienta principalmente a la construcción de soluciones personalizadas, reutilizando componentes, patrones de arquitectura, servicios de integración, mecanismos de seguridad y prácticas de ingeniería desarrolladas en proyectos anteriores.

\section{1.2. Estructura Organizacional}\label{estructura-organizacional}

Para la ejecución del proyecto, Only Simple Solutions dispondría de una estructura organizacional de carácter funcional, liderada por un Jefe de Proyecto responsable de la dirección integral del contrato y de la coordinación de los distintos frentes técnicos, operacionales y de gestión.

Para asegurar que ninguna variable del proyecto quede al azar, hemos estructurado al equipo no como departamentos aislados, sino como seis frentes de control simultáneos. Mientras Arquitectura y Datos aseguran la solidez técnica de la solución, Seguridad y Calidad actúan como filtros independientes para blindar el producto antes de su salida. Finalmente, Operación y Gestión del Cambio garantizan que la tecnología aterrice suavemente en el negocio y sea adoptada por los usuarios. Cada uno de estos frentes está dirigido por un líder especialista, asegurando que el Jefe de Proyecto cuente con garantes expertos para cada riesgo crítico.

\begin{figure}[htbp]
\centering
\includegraphics[width=\linewidth]{figuras/image1-08a792bcec.png}
\caption{Organigrama Only Simple Solutions}
\label{fig:organigrama-only-simple-solutions}
\end{figure}

La dotación total asignada al proyecto estaría compuesta por 16 profesionales, distribuidos de la siguiente manera:

\begin{itemize}
\tightlist
\item
  1 Jefe de Proyecto: responsable de la dirección integral del proyecto, supervisando plazos, recursos, riesgos, avances, coordinación entre áreas y cumplimiento de los objetivos y compromisos establecidos.
\item
  1 Arquitecto de Solución y 2 Ingenieros de Solución/Desarrollo: responsables del diseño técnico de la solución, definición de componentes e integraciones, desarrollo de los componentes requeridos y aplicación de los lineamientos arquitectónicos definidos para el proyecto.
\item
  1 Encargado de Seguridad de la Información y 1 Analista de Seguridad: responsables de definir, implementar y supervisar las medidas de seguridad aplicables a la solución, incluyendo protección de la información, control de accesos, gestión de riesgos y seguimiento de controles de seguridad.
\item
  1 Líder de Datos y 2 Ingenieros de Datos: responsables de la arquitectura, integración, transformación, calidad y migración de los datos asociados al proyecto, asegurando su disponibilidad, consistencia e integridad.
\item
  1 Líder de Calidad y 2 Ingenieros QA: responsables de planificar y ejecutar las actividades de aseguramiento de calidad, pruebas funcionales y técnicas, registro y seguimiento de incidencias y verificación del cumplimiento de los requerimientos definidos.
\item
  1 Líder de Operación y 1 Ingeniero DevOps/Soporte: responsables de preparar y coordinar la operación de la solución, incluyendo monitoreo, continuidad, soporte, gestión de incidentes y automatización de los procesos necesarios para su funcionamiento.
\item
  1 Líder de Implantación y Gestión del Cambio y 1 Especialista de Capacitación e Implantación: responsables de coordinar la puesta en marcha de la solución, acompañamiento a los usuarios, capacitación y actividades de gestión del cambio requeridas para la adopción de la nueva solución.
\end{itemize}

Esta dotación corresponde al equipo de liderazgo y especialistas responsables de los principales frentes del proyecto. Su estructura permite que cada ámbito crítico cuente con una contraparte claramente identificada, manteniendo al Jefe de Proyecto como punto central de coordinación y escalamiento.

\section{1.3. Gobierno interno Calidad, Seguridad y Conocimiento}\label{gobierno-interno-calidad-seguridad-y-conocimiento}

La construcción y modernización de ecosistemas tecnológicos requiere mucho más que la excelencia técnica. Exige un enfoque integral de gestión y control, capaz de abordar de manera sistemática los riesgos asociados a la seguridad, la calidad, la continuidad operacional y el cumplimiento normativo.

Nuestro modelo se sustenta en tres sistemas de gestión con marco formal, todos derivados de normas internacionales certificables: el sistema de gestión de la calidad conforme a ISO 9001, el sistema de gestión de la seguridad de la información conforme a ISO/IEC 27001 e ISO/IEC 27002, y el sistema de gestión de la continuidad del negocio conforme a ISO 22301. La gestión del conocimiento se apoya en las mismas instancias y se sustenta en la trazabilidad documental exigida por el contrato.

La estructura del gobierno es de tres capas. La primera capa son las políticas, aprobadas por la Gerencia General y vigentes con independencia del contrato. La segunda capa son las instancias: comités y sesiones formales con participación, periodicidad y entregables definidos. La tercera capa son los responsables, con roles nombrados y atribuciones explícitas. Ningún control se aplica en un proyecto si antes no está definido quién lo ejecuta, cada cuánto y con qué evidencia.

\subsection{1.3.1. Gobierno interno de la calidad}\label{gobierno-interno-de-la-calidad}

La política de calidad de Only Simple Solutions declara que la compañía se obliga a entregar soluciones que satisfagan los requisitos acordados con el mandante, a mantener veraz y vigente la información de gestión y a mejorar de manera continua la eficacia de sus procesos. La política está aprobada por la Gerencia General, es conocida por toda la dotación y se revisa en la revisión por la dirección.

\begin{itemize}
\tightlist
\item
  Comité de Calidad y Mejora Continua: Se reúne trimestralmente y es responsable de priorizar las acciones correctivas y preventivas abiertas, revisar los indicadores de calidad de la compañía y aprobar las mejoras de proceso que afecten a más de un proyecto. Lo preside el Líder de Calidad y sesionan los directores de las áreas técnicas.
\item
  Revisión por la dirección: Se realiza anualmente y tiene por objeto comprobar que la política de calidad sigue siendo pertinente, que los objetivos de calidad se alcanzan y que los recursos siguen siendo suficientes.
\item
  Auditorías internas: El programa de auditorías internas cubre cada año los procesos del sistema de gestión. Los hallazgos se registran como no conformidades, con responsable, plazo y evidencia de cierre.
\item
  Junta de no conformidades del proyecto: Es la instancia de primer nivel dentro de cada contrato. La preside el Líder de Calidad del proyecto y sesiona semanalmente durante la construcción y mensualmente durante la operación. Revisa las desviaciones detectadas en las pruebas, los incumplimientos de especificación y las observaciones de auditoría.
\end{itemize}

El Líder de Calidad es responsable del sistema y de la independencia de la función respecto de la ejecución, este aplica el sistema dentro del contrato: planifica las pruebas, verifica la trazabilidad entre cada requisito y su evidencia de aceptación, controla la configuración de lo entregado y registra toda desviación.

\subsection{1.3.2. Gobierno interno de la seguridad de la información}\label{gobierno-interno-de-la-seguridad-de-la-informaciuxf3n}

La política de seguridad de la información establece que la confidencialidad, la integridad y la disponibilidad de la información de los mandantes se protegen mediante un marco de gestión de riesgos, controles de acceso con privilegios mínimos, trazabilidad de las acciones sobre datos sensibles y una gestión de incidentes con respuesta verificable. Se sustenta en ISO/IEC 27001 e ISO/IEC 27002, y se extiende a los servicios en nube con los controles de ISO/IEC 27017 e ISO/IEC 27018. Los proyectos que manejan datos de tarjetas de pago se gestionan además conforme a PCI DSS, y los que manejan información personal conforme a ISO/IEC 27701.

Toda la información de los mandantes se clasifica en cuatro niveles: pública, interna, confidencial y restringida, con reglas explícitas de custodia, transferencia, almacenamiento y destrucción. La información crediticia y la información de tarjetas se clasifican como restringidas.

\begin{itemize}
\tightlist
\item
  Comité de Seguridad de la Información: Se reúne mensualmente. Preside el Encargado de Seguridad de la Información y participan el Líder de Datos, el Líder de Operación y los responsables de seguridad de los proyectos activos. Sus funciones son revisar el registro de riesgos de seguridad, resolver los hallazgos de auditoría, aprobar los cambios que afecten a la clasificación de la información y evaluar los incidentes cerrados.
\item
  Mesa de análisis de incidentes: Se convoca dentro de las primeras cuatro horas de declarado un incidente de seguridad relevante. Participa el responsable de seguridad del proyecto afectado, el Encargado de Seguridad de la Información y el Jefe de Proyecto. Determina la clasificación del incidente, la contención, la comunicación al mandante y la necesidad de notificación.
\item
  Revisión de acceso: Se realiza cada trimestre y cada vez que cambia el equipo de un proyecto. Verifica que los accesos vigentes correspondan a las personas que realmente ejercen la función.
\end{itemize}

El Encargado de Seguridad de la Información responde por la política, el registro de riesgos corporativo y la relación con la certificación, y además aplica los controles: define el modelado de amenazas de la solución, mantiene la matriz de controles trazable a ISO/IEC 27001 e ISO/IEC 27002, ejecuta las pruebas de seguridad y de intrusión, y gestiona la respuesta a incidentes.

\subsection{1.3.3. Gobierno interno de la gestión del conocimiento}\label{gobierno-interno-de-la-gestiuxf3n-del-conocimiento}

La continuidad del proyecto se asegura mediante una distribución adecuada de la documentación y el conocimiento, evitando que la salida de cualquier integrante comprometa la operación. Para ello, la política de la compañía establece que todo conocimiento necesario para operar el sistema debe ser documentado, versionado y transferido, evitando que una sola persona concentre de manera exclusiva el conocimiento de un componente crítico.

\begin{itemize}
\tightlist
\item
  Registro de decisiones: Toda decisión técnica con más de una alternativa se documenta en un registro de decisiones de arquitectura que declara el problema, las alternativas evaluadas, el criterio de selección y las consecuencias aceptadas. Las decisiones del proyecto se vinculan con las decisiones de la solución.
\item
  Documentación bajo control: Los manuales de operación, las guías de resolución de fallas, los libros de operación y los procedimientos de conmutación se mantienen en el repositorio del proyecto, con versión, fecha y responsable. Se entregan al mandante en formato editable y en propiedad de este.
\item
  Base de conocimiento: La Mesa de Servicio mantiene una base de conocimiento de incidentes, problemas y soluciones que se enriquece con cada caso resuelto. La base se revisa mensualmente y reporta su uso y su utilidad percibida.
\item
  Transferencia en doble dirección: Cada conocimiento crítico tiene al menos dos personas capaces de aplicarlo. La transferencia se realiza con acompañamiento formal entre pares durante los primeros meses de cada frente y con una sesión de traspaso antes del cierre de cada etapa.
\item
  Comunidades de práctica: Las distintas áreas especializadas de la compañía sostienen reuniones técnicas periódicas sobre materias comunes, cuyo resultado se documenta y se comparte con todos los proyectos.
\end{itemize}

El gobierno de la gestión del conocimiento se ejerce trimestralmente revisando el estado de la base de conocimiento y de los registros de decisión. El Jefe de Proyectos es responsable de que la distribución del conocimiento se mantenga en cada contrato, y el Líder de cada frente es responsable de que la documentación de su frente esté completa y actualizada al cierre de cada hito.

\section{1.4. Experiencia y Certificaciones}\label{experiencia-y-certificaciones}

Only Simple Solutions cuenta con experiencia en soluciones para el comercio minorista, los servicios financieros y el comercio electrónico. Los tres proyectos de referencia, ejecutados como contratista principal, han permitido desarrollar las siguientes capacidades:

\begin{itemize}
\tightlist
\item
  Continuidad operacional y arquitectura híbrida. La modernización de puntos de venta con funcionamiento offline aportó experiencia en la integración de servidores locales con servicios de nube y en la continuidad de operaciones ante interrupciones de conectividad. Esta experiencia fortalece nuestra capacidad para diseñar soluciones distribuidas y gestionar su sincronización.
\item
  Gestión de información financiera. El desarrollo de una plataforma de consentimiento seguro aportó experiencia en el resguardo de información sensible y la trazabilidad de operaciones. Estos conocimientos fortalecen nuestra capacidad para incorporar controles de seguridad y mecanismos de consulta de evidencias en las soluciones.
\item
  Integración omnicanal y alta demanda. La implementación de un motor ATP e integración de marketplace aportó experiencia en la coordinación de inventario y pedidos entre tiendas físicas, comercio electrónico y vendedores externos. Esta experiencia fortalece nuestra capacidad para integrar sistemas y dimensionar plataformas sujetas a variaciones de carga.
\end{itemize}

En los tres proyectos, especificados en el Anexo A.2, Only Simple Solutions participó como contratista principal. En conjunto, estos antecedentes muestran experiencia en soluciones distribuidas que combinan continuidad operacional, gestión de información sensible e integración entre sistemas.

\subsection{1.4.1 Certificaciones}\label{certificaciones}

En Only Simple Solutions, la protección de la información de nuestros clientes y la continuidad operativa son pilares fundamentales de nuestra propuesta de valor. Para garantizar los más altos estándares en la ejecución de nuestros servicios, respaldamos nuestra gestión bajo un marco de certificaciones internacionales:

\begin{itemize}
\tightlist
\item
  ISO/IEC 27001 (Sistema de Gestión de Seguridad de la Información): Nos permite establecer un marco robusto para identificar, evaluar y tratar los riesgos que afectan la confidencialidad, integridad y disponibilidad de los datos. En los proyectos, esto se traduce en estrictos controles de acceso, segregación de responsabilidades, gestión de identidad y monitoreo continuo.
\item
  ISO/IEC 27701 (Gestión de la Privacidad de la Información - PIMS): Nos permite establecer responsabilidades y controles sobre el tratamiento de datos personales, definiendo quién puede utilizarlos, con qué finalidad y bajo qué condiciones. La adoptamos porque el cuidado de la privacidad requiere procedimientos que permitan revisar y demostrar cómo se utiliza la información que nuestros clientes nos confían.
\item
  ISO 9001 (Sistema de Gestión de Calidad): Nos permite trabajar con requisitos acordados, verificar nuestras entregas y corregir las desviaciones detectadas. Su valor está en mantener una forma consistente de trabajar, con controles y evidencias que permitan comprobar que lo entregado responde a lo solicitado, y mejorar nuestros procesos para reducir la repetición de errores.
\item
  ISO 22301 (Sistema de Gestión de la Continuidad del Negocio): Nos otorga la capacidad de respuesta y resiliencia necesaria para mantener la operatividad de los servicios ante cualquier imprevisto o interrupción, garantizando la estabilidad que nuestros clientes requieren. La adoptamos porque la continuidad requiere anticipar escenarios de falla, definir responsables y coordinar la recuperación, reduciendo el impacto sobre los servicios que prestamos y las operaciones de nuestros clientes.
\end{itemize}

Elegimos estas normas por su reconocimiento internacional y porque ofrecen criterios comunes para evaluar la seguridad, la privacidad, la calidad y la continuidad del negocio. Su certificación permite que organismos independientes revisen nuestros procesos, entregando a nuestros clientes un respaldo verificable al momento de confiar en nuestros servicios.

\section{1.5. Estructura para Proyecto}\label{estructura-para-proyecto}

La empresa organizará la ejecución del proyecto de Multitiendas Ancoa S.A. mediante una jefatura central y frentes especializados de arquitectura, desarrollo, datos, seguridad, calidad y operación. Esta estructura permitirá coordinar la implementación en las 22 tiendas y los dos centros de distribución, atender las necesidades del canal digital y del negocio financiero, y mantener la continuidad de los servicios durante la transición.

La organización cubrirá los 56 meses del contrato y adaptará su participación a las etapas de desarrollo, marcha blanca y operación. Su diseño considera la coexistencia de sistemas nuevos y legados, la separación entre los ámbitos de retail y crédito, la migración de una cartera financiera activa y las restricciones del calendario comercial de Ancoa.

Para coordinar el trabajo, se utilizará Jira como herramienta de registro y seguimiento de tareas, responsables, plazos e incidencias, permitiendo visualizar el avance e identificar oportunamente los pendientes que puedan afectar las entregas. Discord se utilizará para la comunicación cotidiana del equipo y la coordinación entre los distintos frentes. Los acuerdos que afecten tareas, plazos o entregables quedarán registrados en Jira para mantener su trazabilidad. Estas herramientas complementarán los roles funcionales del proyecto y las instancias de coordinación con el mandante descritas a continuación.

La estructura del proyecto se conformará con los mismos roles presentados en el punto 1.2 de este subdocumento, cuyas responsabilidades se aplicarán a las necesidades específicas del proyecto de Multitiendas Ancoa S.A.:

\begin{itemize}
\tightlist
\item
  1 Jefe de Proyecto: Dirigirá la ejecución, coordinará los frentes, controlará avances y recursos, gestionará riesgos y tramitará las decisiones y solicitudes de cambio ante el mandante.
\item
  1 Arquitecto de Solución: Definirá y mantendrá la coherencia de la arquitectura híbrida. Coordinará el levantamiento de las nueve plataformas y catorce interfaces existentes, las integraciones y la separación técnica entre retail y negocio financiero.
\item
  1 Encargado de Seguridad de la Información: Establecerá y verificará los controles de acceso, segregación, protección de datos, trazabilidad y respuesta a incidentes. Coordinará con las contrapartes de Ancoa la revisión de los intercambios de información entre ambos negocios.
\item
  1 Líder de Datos: Dirigirá el análisis, saneamiento, migración y conciliación de la información. Coordinará la definición de fuentes de verdad y la migración de la cartera de crédito, con controles sobre saldos, cuotas, repactaciones y evidencias asociadas.
\item
  1 Líder de Calidad: Coordinará la verificación de requisitos, las pruebas y las evidencias de aceptación. Mantendrá independencia funcional respecto de la construcción de los componentes evaluados y reportará las desviaciones detectadas antes de cada liberación.
\item
  1 Líder de Operación: Preparará y conducirá la operación de la plataforma, la mesa de servicio, el monitoreo y los procedimientos de continuidad y recuperación. Coordinará el soporte en terreno y la transferencia de conocimiento al equipo técnico de Ancoa.
\item
  1 Líder de Implantación y Gestión del Cambio: Coordinará el despliegue de la solución y el acompañamiento durante las marchas blancas, considerando la continuidad operacional y las ventanas de congelamiento comercial. Dirigirá la capacitación por perfil y turno y gestionará la adopción, considerando la rotación del personal y las incorporaciones de temporada.
\end{itemize}

Coordinación y gobierno con el mandante

Comité Ejecutivo. Se reunirá mensualmente con participación de la gerencia de la empresa, el patrocinador de Ancoa y el Administrador del Contrato. Abordará decisiones estratégicas, asuntos escalados, aprobación de cambios de alcance y revisión de riesgos mayores.

Comité de Proyecto. Se reunirá quincenalmente y estará integrado por el Jefe de Proyecto y la Contraparte Técnica de Ancoa. Revisará el avance del cronograma, el estado de los entregables, las desviaciones y las decisiones operativas.

Comité de Arquitectura. Se reunirá mensualmente con participación del Arquitecto de Solución, el Encargado de Seguridad de la Información y los referentes técnicos de Ancoa. Se encargará de aprobar las decisiones de arquitectura y revisar la deuda técnica y los riesgos técnicos.

Comité de Operación. Se reunirá mensualmente desde el mes 13, con participación del Líder de Operación, la mesa de servicio y la Contraparte Técnica de Ancoa. Revisará el cumplimiento de los niveles de servicio, los incidentes, los problemas y el plan de mejora continua.

Reunión de seguimiento. Se realizará semanalmente con los equipos de trabajo de ambas partes. Permitirá coordinar las actividades, gestionar dependencias e impedimentos y establecer los compromisos de la semana.

\section{1.6. Alianzas}\label{alianzas}

Only Simple Solutions articula sus alianzas tecnológicas en torno a tres ámbitos complementarios: servicios de nube, seguridad de redes e infraestructura local. Las relaciones con Microsoft, Fortinet y Dell Technologies, directamente o mediante integradores autorizados según el acuerdo correspondiente, complementan nuestras capacidades de desarrollo e integración de software y respaldan el enfoque de soluciones híbridas.

La alianza con Microsoft se vincula con el desarrollo e integración de soluciones sobre Azure. Su aporte se concentra en la infraestructura de nube y los servicios utilizados para desplegar, administrar y monitorear aplicaciones y datos.

La alianza con Fortinet se orienta a la protección de las comunicaciones y la seguridad de las redes. Su ámbito comprende soluciones de cortafuegos, segmentación y conectividad segura entre instalaciones locales y servicios de nube. Estas capacidades complementan el desarrollo de software al abordar la protección de los entornos donde se ejecutan las aplicaciones y circula la información.

La alianza con Dell Technologies, directamente o mediante un integrador de su infraestructura, se relaciona con los servidores y componentes físicos necesarios para las soluciones que requieren procesamiento local.

Estas relaciones complementan las capacidades de Only Simple Solutions para integrar aplicaciones, comunicaciones e infraestructura.

\section*{Referencias}
\addcontentsline{toc}{section}{Referencias}

\section*{Declaración de uso de IA}
\addcontentsline{toc}{section}{Declaración de uso de IA}

El equipo declara que el uso de herramientas de inteligencia artificial generativa en este subdocumento se limitó a las actividades indicadas en la siguiente tabla. Todo contenido, dato, cálculo, decisión de diseño, diagrama, conclusión y referencia fue revisado y verificado por las personas identificadas, quienes asumen como propio el resultado final y la responsabilidad íntegra sobre su exactitud, originalidad y coherencia técnica.

\begin{longtable}[]{@{}
  >{\raggedright\arraybackslash}p{(\columnwidth - 10\tabcolsep) * \real{0.1667}}
  >{\raggedright\arraybackslash}p{(\columnwidth - 10\tabcolsep) * \real{0.1667}}
  >{\raggedright\arraybackslash}p{(\columnwidth - 10\tabcolsep) * \real{0.1667}}
  >{\raggedright\arraybackslash}p{(\columnwidth - 10\tabcolsep) * \real{0.1667}}
  >{\raggedright\arraybackslash}p{(\columnwidth - 10\tabcolsep) * \real{0.1667}}
  >{\raggedright\arraybackslash}p{(\columnwidth - 10\tabcolsep) * \real{0.1667}}@{}}
\caption{Declaración de uso de IA}\label{tab:declaracion-uso-ia}\\
\toprule\noalign{}
\begin{minipage}[b]{\linewidth}\raggedright
Sección
\end{minipage} & \begin{minipage}[b]{\linewidth}\raggedright
Herramienta
\end{minipage} & \begin{minipage}[b]{\linewidth}\raggedright
Finalidad de uso
\end{minipage} & \begin{minipage}[b]{\linewidth}\raggedright
Nivel en texto
\end{minipage} & \begin{minipage}[b]{\linewidth}\raggedright
Nivel en diagramas
\end{minipage} & \begin{minipage}[b]{\linewidth}\raggedright
Revisión humana
\end{minipage} \\
\midrule\noalign{}
\endfirsthead
\toprule\noalign{}
\begin{minipage}[b]{\linewidth}\raggedright
Sección
\end{minipage} & \begin{minipage}[b]{\linewidth}\raggedright
Herramienta
\end{minipage} & \begin{minipage}[b]{\linewidth}\raggedright
Finalidad de uso
\end{minipage} & \begin{minipage}[b]{\linewidth}\raggedright
Nivel en texto
\end{minipage} & \begin{minipage}[b]{\linewidth}\raggedright
Nivel en diagramas
\end{minipage} & \begin{minipage}[b]{\linewidth}\raggedright
Revisión humana
\end{minipage} \\
\midrule\noalign{}
\endhead
\bottomrule\noalign{}
\endlastfoot
1.1. & OpenCode Big Pickle & Creación de datos ficticios de la compañía. & Alto & Ninguno & Eliseo Guarda \\
1.1.1. & Codex GPT 5.6 Sol & Borrador de lista de capacidades útiles para el caso. & Medio & Ninguno & Alex Alfaro \\
1.1.2. & Codex GPT 5.6 Sol & Borrador de lista de servicios útiles para el caso. & Medio & Ninguno & Alex Alfaro \\
1.2. & OpenCode Big Pickle & Borrador de lista de cargos. Creación de la Figura 1. & Medio & Alto & Eliseo Guarda \\
1.3. & OpenCode Big Pickle & Resumen de las normas ISO a seguir; borrador de la sección. & Medio & Ninguno & Eliseo Guarda \\
1.3.1. & OpenCode Big Pickle & Ideación y redacción del gobierno interno de la compañía. & Alto & Ninguno & Eliseo Guarda \\
1.3.2. & OpenCode Big Pickle & Ideación y redacción del gobierno interno de la compañía. & Alto & Ninguno & Eliseo Guarda \\
1.3.3. & OpenCode Big Pickle & Ideación y redacción del gobierno interno de la compañía. & Alto & Ninguno & Eliseo Guarda \\
1.4. & Codex GPT 5.6 Sol & Resumen de proyectos anteriormente ideados. & Medio & Ninguno & Alex Alfaro \\
1.4.1. & Codex GPT 5.6 Sol & Borrador de lista de posibles certificaciones útiles para el caso. & Medio & Ninguno & Alex Alfaro \\
1.5. & Codex GPT 5.6 Sol & Redacción de la sección. & Alto & Ninguno & Alex Alfaro \\
1.6. & Codex GPT 5.6 Sol & Redacción de la sección. & Alto & Ninguno & Alex Alfaro \\
Anexo A.1 & Codex GPT 5.6 Sol & Fue utilizado para la creación de la página web & Ninguno & Ninguno & Manuel Aguilera \\
Formulario T-6 & Codex GPT 5.6 Sol & Creación de proyectos ficticios de la compañía. & Alto & Ninguno & Alex Alfaro \\
\end{longtable}

A continuación, se detalla la escala de niveles para completar la tabla anterior.

\begin{longtable}[]{@{}
  >{\raggedright\arraybackslash}p{(\columnwidth - 4\tabcolsep) * \real{0.3333}}
  >{\raggedright\arraybackslash}p{(\columnwidth - 4\tabcolsep) * \real{0.3333}}
  >{\raggedright\arraybackslash}p{(\columnwidth - 4\tabcolsep) * \real{0.3333}}@{}}
\caption{Niveles de uso de IA}\label{tab:niveles-uso-ia}\\
\toprule\noalign{}
\begin{minipage}[b]{\linewidth}\raggedright
Nivel
\end{minipage} & \begin{minipage}[b]{\linewidth}\raggedright
Texto
\end{minipage} & \begin{minipage}[b]{\linewidth}\raggedright
Diagramas
\end{minipage} \\
\midrule\noalign{}
\endhead
\bottomrule\noalign{}
\endlastfoot
Ninguno & Sin uso de IA. & Sin uso de IA. \\
Bajo & Corrección ortográfica, de estilo o reformulación de frases escritas por el grupo. & Sugerencias de formato o disposición sobre un diagrama elaborado por el grupo. \\
Medio & Borradores o síntesis de partes que el grupo reescribió y verificó. & Diagrama generado con asistencia (por ejemplo, código Mermaid o PlantUML) a partir de un modelo definido por el grupo. \\
Alto & Texto generado sustancialmente por IA y editado por el grupo. & Diagrama generado sustancialmente por IA. \\
\end{longtable}

\ossFinalPage

\end{document}
